\documentclass[a4paper,11pt]{article}
\usepackage[utf8]{inputenc}
\usepackage[T1]{fontenc}
\usepackage[french]{babel}
\usepackage[left=1cm,right=1cm,top=.5cm,bottom=0cm]{geometry}
\usepackage{enumitem}
\usepackage{hyperref}
\usepackage{xcolor}
\usepackage{titlesec}
\usepackage{fontawesome5}
\usepackage{lmodern}
\usepackage[normalem]{ulem}

% --- Couleurs Schneider Electric ---
\definecolor{primary}{RGB}{63, 150, 60}
\definecolor{secondary}{RGB}{80, 80, 80}

% --- Config Liens ---
\hypersetup{colorlinks=true, linkcolor=primary, filecolor=primary, urlcolor=primary}

% --- Titres ---
\titleformat{\section}{\large\bfseries\color{primary}\uppercase}{}{0em}{}[\titlerule]
\titlespacing{\section}{0pt}{8pt}{5pt}

% --- Commandes ---
\newcommand{\resumeItem}[1]{\item\small{#1}}

\begin{document}

% --- EN-TÊTE ---
\begin{center}
    {\Large \textbf{Loïc Aron MBASSI EWOLO}} \\ \vspace{4pt}
    \textbf{\large Alternance Développeur Logiciel Embarqué Réseaux Industriels} \\
    \textit{\small Disponible dès septembre 2026 pour 12 mois} \\ \vspace{4pt}
    \small
    \faMapMarker\ Cergy, France (Mobile France) \quad
    \faPhone\ (+33) 7 59 52 33 98 \quad
    {\color{primary}\faEnvelope\ \href{mailto:loicaron.mbassiewolo@ensea.fr}{loicaron.mbassiewolo@ensea.fr}} \\ \vspace{2pt}
    {\color{primary}\faLinkedin\ \href{https://www.linkedin.com/in/loic-mbassi/}{linkedin.com/in/loic-mbassi}} \quad
    {\color{primary}\faGithub\ \href{https://github.com/Nameless0l}{github.com/Nameless0l}} \quad
    {\color{primary}\faGlobe\ \href{https://mbassiloic.tech}{mbassiloic.tech}}
\end{center}

\vspace{-6pt}

% --- PROFIL ---
\noindent\small{Développement embarqué C/C++ sur Nvidia Jetson avec intégration capteurs réseau (Lidar 3D/UDP, caméra de profondeur) et programmation système C sous Linux (IPC, mémoire partagée, signaux POSIX). Double diplôme ENSEA/Polytechnique Yaoundé en systèmes embarqués, signal et IA. Prototypage IoT sur STM32 (capteurs, fusion données) et optimisation de modèles ML pour microcontrôleurs (TinyML, TF Lite). Expérience en architecture microservices avec communication asynchrone (RabbitMQ).}

% --- FORMATION ---
\section{Formation}
\begin{itemize}[leftmargin=*, label=\tiny$\bullet$, nosep]
    \item \textbf{ENSEA -- École Nationale Supérieure de l'Électronique et de ses Applications} \hfill \textit{2025 -- 2027} \\
    \small{Double diplôme d'Ingénieur -- Systèmes Embarqués, FPGA/VHDL, Traitement du Signal, Réseaux, Deep Learning.}
    \item \textbf{École Polytechnique de Yaoundé (1\textsuperscript{ère} école d'ingénieur CEMAC)} \hfill \textit{2021 -- 2025} \\
    \small{Diplôme d'Ingénieur de Conception (Master 2) : \textbf{Mathématiques Appliquées}, Génie Logiciel, Algorithmique.}
\end{itemize}

% --- COMPÉTENCES ---
\section{Compétences Techniques}
\begin{itemize}[leftmargin=*, nosep]
    \item \small{\textbf{Embarqué \& RTOS :} C embarqué, C++, STM32, Arduino, Raspberry Pi, Nvidia Jetson, contraintes temps réel, FPGA/VHDL.}
    \item \small{\textbf{Réseaux \& Protocoles :} Couches basses Ethernet, communication inter-processus (IPC, SHM), RabbitMQ, API REST, UDP.}
    \item \small{\textbf{Outils \& Debug :} Docker, Git, CI/CD, compilation croisée, Python (scripting, tests), MATLAB/Simulink.}
    \item \small{\textbf{Capteurs \& Intégration :} Lidar 3D (VLP16), caméra de profondeur, IMU, flex sensors, SLAM, fusion de capteurs.}
    \item \small{\textbf{Langues :} Français (natif), Anglais (professionnel/technique -- rédaction scientifique INRIA).}
\end{itemize}

% --- EXPÉRIENCES / PROJETS ---
\section{Expériences \& Projets Embarqués}
\begin{itemize}[leftmargin=*, label={-}]

\item \textbf{Challenge CoHoMa -- Robotique Autonome (ENSEA / Armée française)} \hfill \textit{2025 -- En cours} \\
\textit{Développement embarqué C/C++, intégration capteurs réseau} \hfill \textit{C/C++, Nvidia Jetson, Lidar 3D, Docker}
\vspace{-2pt}
    \begin{itemize}[leftmargin=0.4cm, nosep, label=\tiny$\bullet$]
        \item \small{Développement C/C++ sur Jetson : intégration Lidar VLP16 (flux UDP), caméra de profondeur, SLAM temps réel.}
        \item \small{Optimisation d'inférence YOLOv10 sur GPU embarqué, fusion de capteurs multi-protocoles.}
        \item \small{Containerisation Docker pour reproductibilité et intégration continue multi-plateforme.}
    \end{itemize}
\vspace{3pt}

\item \textbf{Simulation Flotte de Taxis -- Programmation Système C/Linux} \hfill \textit{2024} \\
\textit{Projet académique -- ENSPY} \hfill \textit{C, IPC (SHM, Signaux POSIX), OpenGL, Linux}
\vspace{-2pt}
    \begin{itemize}[leftmargin=0.4cm, nosep, label=\tiny$\bullet$]
        \item \small{Simulation multiprocessus : mémoire partagée, signaux POSIX, ordonnanceur FIFO, communication inter-processus.}
        \item \small{Gestion manuelle de la mémoire, synchronisation de processus concurrents, programmation bas niveau.}
    \end{itemize}
\vspace{3pt}

\item \textbf{Harmony Gloves -- Gants IoT (STM32, Capteurs)} \hfill \textit{2024} \\
\textit{Labo IOTA, École Polytechnique de Yaoundé} \hfill \textit{STM32, Arduino, C, Capteurs IMU/Flex}
\vspace{-2pt}
    \begin{itemize}[leftmargin=0.4cm, nosep, label=\tiny$\bullet$]
        \item \small{Participation à la conception de gants instrumentés : lecture capteurs inertiels, fusion données capteurs + vision.}
        \item \small{Intégration hardware/software, soudure composants, modèle ML embarqué pour classification temps réel.}
    \end{itemize}
\vspace{3pt}

\item \textbf{Stage Recherche IoT -- Edge Computing \& TinyML} \hfill \textit{Oct 2023 -- Jan 2024} \\
\textit{IoT Lab, École Polytechnique de Yaoundé} \hfill \textit{Python, TF Lite, Arduino, Raspberry Pi, C}
\vspace{-2pt}
    \begin{itemize}[leftmargin=0.4cm, nosep, label=\tiny$\bullet$]
        \item \small{Optimisation de modèles ML pour microcontrôleurs : quantification, pruning, contraintes mémoire et temps réel.}
        \item \small{Déploiement TensorFlow Lite sur Arduino/Raspberry Pi, benchmarking de performances embarquées.}
    \end{itemize}
\vspace{3pt}

\item \textbf{Architecture Microservices -- GoFind (Communication Asynchrone)} \hfill \textit{2024} \\
\textit{Projet académique} \hfill \textit{Laravel, NestJS, MongoDB, RabbitMQ, Docker}
\vspace{-2pt}
    \begin{itemize}[leftmargin=0.4cm, nosep, label=\tiny$\bullet$]
        \item \small{Architecture microservices : API Gateway, communication asynchrone (RabbitMQ), protocoles de messagerie.}
        \item \small{Cartographie des flux de données entre services, tests d'intégration, containerisation Docker.}
    \end{itemize}
\vspace{3pt}

\end{itemize}

% --- DISTINCTIONS ---
\section{Distinctions \& Compléments}
\begin{itemize}[leftmargin=*, label=\tiny$\bullet$, nosep]
    \item \small{\textbf{Stage INRIA Saclay} (Équipe TRIBE, avr--août 2026) : pipeline Python, analyse statique, rédaction scientifique anglais.}
    \item \small{\textbf{1\textsuperscript{er} Prix} CITS National 2024 (75 équipes). \textbf{Lead AI Cell} (ENSPY) : workshops embarqué, ML.}
    \item \small{\textbf{Certif.} Supervised ML (Stanford/DeepLearning.AI), Python for Data Science (IBM/Coursera).}
\end{itemize}

\end{document}
